\clearpage\pagestyle{otMain}\setlength{\otBlocEspacement}{2.5mm}\fontsize{11}{15}\selectfont\corrigesChapitre{9}{Développer, réduire et démontrer}
\begin{blocLivre}{cor-9-1}\begin{corrige}{1}{Développer une somme}\demarche{Distribuer sur tous les termes.}
On multiplie chacun des deux termes par 3 : $\res{3x+12}$.
\end{corrige}\end{blocLivre}
\begin{blocLivre}{cor-9-2}\begin{corrige}{2}{Développer une différence}\demarche{Distribuer sur la différence.}
$5\times x-5\times2=\res{5x-10}$. Le signe de la différence est conservé.
\end{corrige}\end{blocLivre}
\begin{blocLivre}{cor-9-3}\begin{corrige}{3}{Réduire des termes semblables}\demarche{Regrouper les multiples de la lettre.}
$7x+2x=9x$. L’expression réduite est $\res{9x+4}$.
\end{corrige}\end{blocLivre}
\begin{blocLivre}{cor-9-4}\begin{corrige}{4}{Un facteur commun}\demarche{Repérer un multiplicateur commun.}
Les deux termes contiennent le facteur 4 : $\res{4(x+5)}$. Développer permet de contrôler.
\end{corrige}\end{blocLivre}
\begin{blocLivre}{cor-9-5}\begin{corrige}{5}{Avec des décimaux}\demarche{Additionner les coefficients.}
$1{,}5+2=3{,}5$, donc $\res{3{,}5x+3}$.
\end{corrige}\end{blocLivre}
\begin{blocLivre}{cor-9-6}\begin{corrige}{6}{Développer puis réduire}\demarche{Procéder en deux étapes.}
$2x+10+3x=\res{5x+10}$. On ne regroupe qu’après avoir distribué.
\end{corrige}\end{blocLivre}
\begin{blocLivre}{cor-9-7}\begin{corrige}{7}{Une différence de termes}\demarche{Regrouper séparément les types de termes.}
$8x-3x=5x$ et $7+2=9$. On obtient $\res{5x+9}$.
\end{corrige}\end{blocLivre}
\begin{blocLivre}{cor-9-8}\begin{corrige}{8}{Factoriser une différence}\demarche{Chercher un facteur dans les deux termes.}
Le facteur commun 7 donne $\res{7(x-3)}$. Le contrôle par développement redonne $7x-21$.
\end{corrige}\end{blocLivre}
\begin{blocLivre}{cor-9-9}\begin{corrige}{9}{Comparer deux programmes}\demarche{Utiliser une identité, au-delà des tests.}
La distributivité donne $4(x+2)=4x+8$. Les deux programmes sont $\res{\text{équivalents}}$ pour toute valeur de $x$.
\end{corrige}\end{blocLivre}
\begin{blocLivre}{cor-9-10}\begin{corrige}{10}{Une erreur à expliquer}\demarche{Tester une valeur qui distingue les écritures.}
Pour $x=2$, $2x+3=7$ mais $5x=10$. L’affirmation est fausse. $\res{2x+3}$ ne se réduit pas davantage : les termes sont de natures différentes.
\end{corrige}\end{blocLivre}
\begin{blocLivre}{cor-9-11}\begin{corrige}{11}{Un périmètre général}\demarche{Revenir à la somme des quatre côtés.}
$P=2(x+4)+2(x+1)=2x+8+2x+2=\res{4x+10}$ cm. Les deux dimensions apparaissent deux fois.
\label{fig-cor-9-11}\figureGeo[corrige]{dist-11}
\end{corrige}\end{blocLivre}
\begin{blocLivre}{cor-9-12}\begin{corrige}{12}{Une aire découpée}\demarche{Calculer une même aire de deux façons.}
Le grand rectangle a pour aire $3(x+5)$. Les parties ont pour aires $3x$ et 15. En additionnant : $\res{3(x+5)=3x+15}$ cm$^2$. Le découpage explique la distributivité.
\label{fig-cor-9-12}\figureGeo[corrige]{dist-12}
\end{corrige}\end{blocLivre}
\begin{blocLivre}{cor-9-13}\begin{corrige}{13}{Trois entiers consécutifs}\demarche{Réduire puis factoriser.}
La somme vaut $n+(n+1)+(n+2)=3n+3=3(n+1)$. C’est $\res{\text{trois fois le nombre du milieu}}$, pour tout entier $n$.
\end{corrige}\end{blocLivre}
\begin{blocLivre}{cor-9-14}\begin{corrige}{14}{Un programme invariant}\demarche{Faire disparaître la lettre par réduction.}
L’expression est $2(x+6)-2x$. On développe : $2x+12-2x=\res{12}$. Les termes dépendant du choix s’annulent, ce qui prouve l’invariance.
\end{corrige}\end{blocLivre}
\begin{blocLivre}{cor-9-15}\begin{corrige}{15}{Une conjecture trompeuse}\demarche{Explorer puis chercher un cas qui peut faire échouer la règle.}
La formule est \texttt{=A2*A2+A2+11}. Pour 1 et 2, on obtient 13 et 17, qui sont premiers : on peut conjecturer la même propriété pour les autres valeurs. Mais pour 10, $100+10+11=\res{121=11\times11}$. Ce contre-exemple réfute l’affirmation générale. Le tableur aide à explorer ; une suite de tests vrais ne suffit pas à prouver.
\end{corrige}\end{blocLivre}
